\chapter{Checklists Acadêmicos}

Os checklists abaixo apoiam a revisão por etapas. Eles não substituem a leitura crítica, a orientação nem os critérios específicos do curso.

\section{Checklist do projeto}

\begin{itemize}
  \item O tema está delimitado e o problema é investigável.
  \item O objetivo geral responde ao problema e os objetivos específicos descrevem etapas necessárias.
  \item A justificativa demonstra relevância acadêmica, social, profissional ou institucional.
  \item O referencial inicial reúne fontes pertinentes e atuais para o objeto.
  \item A metodologia é compatível com os objetivos e descreve análise, limitações e gestão dos dados.
  \item O cronograma é realista e reserva tempo para revisão.
  \item Riscos, autorizações e necessidade de avaliação ética foram verificados antes da coleta.
  \item As referências utilizadas no projeto estão registradas.
\end{itemize}

\section{Checklist para orientação}

\begin{itemize}
  \item A versão enviada está identificada e permite localizar as alterações.
  \item As solicitações anteriores foram atendidas ou justificadas.
  \item As dúvidas para a reunião são objetivas.
  \item Problema, objetivos, método, análise e parte conclusiva permanecem coerentes.
  \item Citações e referências novas estão registradas.
  \item Lacunas, limitações e seções ainda incompletas estão sinalizadas.
  \item O estudante registrou as decisões tomadas durante a orientação.
\end{itemize}

\section{Checklist para banca}

\begin{itemize}
  \item O trabalho possui os elementos obrigatórios aplicáveis.
  \item Resumos, palavras-chave, sumário e listas foram conferidos.
  \item Objetivos, resultados e parte conclusiva são coerentes.
  \item A metodologia permite compreender o percurso realizado.
  \item Figuras, tabelas, quadros e demais objetos são citados e possuem título e fonte.
  \item Citações e referências foram conferidas em ambos os sentidos.
  \item Apêndices e anexos são pertinentes e foram mencionados no texto.
  \item A paginação e a legibilidade foram verificadas no PDF enviado.
  \item O envio foi autorizado conforme a regra do curso.
\end{itemize}

\section{Checklist da versão corrigida}

\begin{itemize}
  \item As correções da banca foram incorporadas ou justificadas à instância competente.
  \item Alterações de conteúdo não criaram contradições em outras seções.
  \item Sumário, listas, paginação e referências foram atualizados.
  \item O PDF abre corretamente e corresponde ao arquivo autorizado para entrega.
  \item Nome, formato, prazo e canal seguem a orientação institucional vigente.
  \item Restrições, embargo, termos ou autorizações foram anexados quando aplicáveis.
\end{itemize}

\section{Checklist de normalização}

\begin{itemize}
  \item Papel, margens, fonte, corpo, espaçamento e parágrafos estão uniformes.
  \item A contagem e a exibição da paginação estão corretas.
  \item A hierarquia de seções é progressiva e não ultrapassa o nível quinário.
  \item Títulos do texto e do sumário coincidem.
  \item Siglas, abreviaturas, alíneas e subalíneas seguem um padrão estável.
  \item Citações diretas possuem localização quando aplicável.
  \item Referências apresentam dados suficientes e não contêm marcadores não preenchidos.
  \item Elementos pré-textuais não aparecem no sumário.
\end{itemize}

\section{Checklist de acessibilidade}

\begin{itemize}
  \item O PDF é pesquisável e possui título e idioma definidos nos metadados.
  \item A ordem de leitura e a hierarquia de títulos são compreensíveis.
  \item Imagens informativas possuem texto alternativo ou descrição equivalente.
  \item Contraste, tamanho do texto e espaçamento favorecem a leitura.
  \item Informações não dependem exclusivamente de cor.
  \item Tabelas possuem cabeçalhos claros e estrutura simples.
  \item Links têm texto descritivo e funcionam no arquivo entregue.
  \item Slides e pôsteres, quando utilizados, seguem os mesmos princípios.
\end{itemize}

\section{Checklist de integridade}

\begin{itemize}
  \item Ideias, textos, dados, imagens e códigos de terceiros estão atribuídos.
  \item Paráfrases foram realmente reescritas e continuam acompanhadas de fonte.
  \item Não há referências inventadas, dados fabricados ou resultados manipulados.
  \item Reutilização de produção anterior foi identificada e autorizada quando necessário.
  \item Contribuições de outras pessoas não foram apresentadas como autoria própria.
  \item O uso de inteligência artificial, quando houver, respeita as regras aplicáveis e foi declarado quando exigido.
  \item Informações pessoais, sigilosas ou protegidas não foram expostas indevidamente.
  \item O estudante revisou o conteúdo e assume responsabilidade pela versão entregue.
\end{itemize}
